\documentclass[10pt,a4paper]{amsart}
\usepackage[margin=19mm]{geometry}
\usepackage{amsmath,amssymb,mathtools}
\usepackage[scr=rsfs,cal=euler]{mathalfa}
\usepackage{xfrac}
\usepackage[unicode,hidelinks,bookmarks=false]{hyperref}
\newcommand{\ie}{i.e.\ }
\newcommand{\cf}{cf.\ }
\newcommand{\resp}{respectively\ }
\newcommand{\Subsection}{\subsection}
\providecommand{\nobreakdash}{\nobreak\hbox{-}}
\setlength{\emergencystretch}{2em}
\multlinegap0pt
\usepackage{bm}
\usepackage{bbm}
\usepackage{dsfont}
\usepackage{xspace}
\usepackage{cancel}
\usepackage{mathtools}
\usepackage{cleveref}
\usepackage{amsthm}
\makeatletter
\@ifundefined{lemma}{\newtheorem{lemma}{Lemma}}{}
\makeatother
\DeclareMathOperator{\dist}{dist}
\DeclarePairedDelimiter{\abs}{\lvert}{\rvert}
\DeclarePairedDelimiter{\norm}{\lVert}{\rVert}
\DeclarePairedDelimiter{\paren}{(}{)}
\newcommand{\CC}[0]{\mathbb{C}}
\newcommand{\RR}[0]{\mathbb{R}}
\newcommand{\hd}[1]{\mathcal{H}^{#1}}
\renewcommand\Re{\operatorname{Re}}%
\renewcommand{\restriction}{\big|}
\renewcommand{\subset}{\subseteq}
\title[The Usual Square Function on Weakly Flat Sets]{The Usual Square Function on Weakly Flat Sets}
\author{Benjamin Jaye, Tobias Wang}
\date{Source: arXiv:2606.24461v1 (CC BY 4.0)}
\begin{document}
\maketitle

\noindent\emph{Source and licence.} The Usual Square Function on Weakly Flat Sets. Version arXiv:2606.24461v1 (CC BY 4.0), distributed under the Creative Commons Attribution 4.0 International licence, as recorded in the arXiv licence metadata for this exact version and shown on the abstract page https://arxiv.org/abs/2606.24461v1. Full rights notice: licenses/arxiv-2606.24461-CC-BY-4.0.txt.\par\smallskip

\noindent\textit{Editorial scope.} Counted unit: the Lemma whose source label is \texttt{eq:segment-usfe-appendix} in the article (source0.tex lines 2741--2749, source label \texttt{eq:segment-usfe-appendix}), delivered together with the article's own complete original proof of it (source0.tex lines 2750--2800, one \texttt{proof} environment of 51 lines). No mathematics is written, repaired or re-derived: statement and proof are the article's own text line for line. Required context retained uncounted: none. Every definition, display and label of those lines is retained unchanged. Omitted from this package and not redistributed: 1--2740, 2801--2964. Nothing inside a retained excerpt is omitted or altered: every retained body is delivered exactly as the article's lines are written, so no omission receipt of any kind accompanies this package. Citations: the retained excerpts carry no citation command at all, so no bibliography entry of the article is transcribed and no reference is redistributed. Reference adaptation: none was needed and none was made; every reference inside the delivered unit resolves inside it, so no unresolved reference or citation is produced. Printed slips kept literal: no printed word, symbol, hypothesis or number of the counted statement or of its proof was altered. Layout and class: the article's own class is \texttt{amsart} with packages including geometry, babel, fontenc, uniinput, csquotes, comment, enumerate, enumitem, amsmath, amssymb, amsthm, thmtools, esint, stmaryrd; the wrapper is the lane's pinned \texttt{amsart} editorial wrapper at 10pt on A4, which restates the theorem-like environments the retained text needs and adds the article's own macro definitions that the retained text needs (\texttt{\textbackslash CC}, \texttt{\textbackslash RR}, \texttt{\textbackslash Re}, \texttt{\textbackslash abs}, \texttt{\textbackslash dist}, \texttt{\textbackslash hd}, \texttt{\textbackslash norm}, \texttt{\textbackslash paren}, \texttt{\textbackslash restriction}, \texttt{\textbackslash subset}), with extra wrapper packages (bm, bbm, dsfont, xspace, cancel, mathtools, cleveref). The delivered numbering is the unit's own. Assets: no retained span contains a figure environment, a graphics-inclusion command, a PDF-page inclusion, a table or a TikZ picture; no image file is copied into this package, the package declares no asset dependency and no image file is redistributed with it. Licence: version arXiv:2606.24461v1 of this article is distributed under the Creative Commons Attribution 4.0 International licence, as recorded in the arXiv licence metadata for this exact version and shown on the abstract page https://arxiv.org/abs/2606.24461v1. A full rights notice is delivered at licenses/arxiv-2606.24461-CC-BY-4.0.txt. Characters: every retained excerpt is pure ASCII, so the delivered engine, XeLaTeX with its default Latin Modern text font, reports no missing character.
\begin{lemma}[Segment square function estimate]
  There is an absolute constant $C$ such that, for every line segment $S\subset\CC$ and every $g\in L^2(\hd1\restriction_S)$,
  \begin{align}\label{eq:segment-usfe-appendix}
    \int_{\CC\setminus S}
    \abs[\Big]{\int_S \frac{g(\xi)}{(z-\xi)^2}\,d\hd1(\xi)}^2
    \dist(z,S)\,dm_2(z)
    \le C\int_S |g(\xi)|^2\,d\hd1(\xi) .
  \end{align}
\end{lemma}

\begin{proof}
  By translation, rotation, and dilation invariance, it is enough to consider $S=[0,L]\subset\RR$.  Extend $g$ by zero to all of $\RR$ and write
  \begin{align}
    Tg(z):=\int_0^L \frac{g(t)}{(z-t)^2}\,dt .
  \end{align}
  We split $\CC\setminus S$ into the vertical strip over the segment and the two endpoint regions.

  First consider the strip $\{z=x+iy:0<x<L,\ y\ne0\}$.  There $\dist(z,S)=|y|$.  For $y>0$, Plancherel's theorem and the standard Fourier transform of $(x+iy)^{-2}$ give
  \begin{align}
    \int_{\RR}|Tg(x+iy)|^2\,dx
    \lesssim
    \int_{\RR}|\xi|^2e^{-c y|\xi|}|\widehat g(\xi)|^2\,d\xi
  \end{align}
  with an absolute constant $c>0$.  Therefore
  \begin{align}
    \int_0^\infty\int_{\RR}|Tg(x+iy)|^2y\,dx\,dy
    &\lesssim
    \int_{\RR}|\widehat g(\xi)|^2|\xi|^2
      \int_0^\infty y e^{-c y|\xi|}\,dy\,d\xi  \\
    &\lesssim \int_{\RR}|\widehat g(\xi)|^2\,d\xi
    \lesssim \int_0^L |g(t)|^2\,dt .
  \end{align}
  The lower half-plane is analogous.
  Restricting the $x$-integral to $0<x<L$ gives the desired estimate in the strip.

  It remains to control the endpoint regions.  We treat $\Omega_-:=\{z:\Re z<0\}$; the region $\Omega_+:=\{z:\Re z>L\}$ is the same after the change of variables $t\mapsto L-t$ and $z\mapsto L-z$.  Write $z=\rho e^{i\theta}$ in $\Omega_-$, so that $\pi/2<\theta<3\pi/2$.  Then $\dist(z,S)=|z|=\rho$, and for $t\ge0$,
  \begin{align}
    |z-t|^2=|z|^2+t^2-2t\Re z\ge \rho^2+t^2 .
  \end{align}
  Hence
  \begin{align}
    \int_{\Omega_-}|Tg(z)|^2\dist(z,S)\,dm_2(z)
    &\lesssim
    \int_0^\infty
      \paren[\Big]{\int_0^L\frac{|g(t)|}{\rho^2+t^2}\,dt}^2\rho^2\,d\rho \\
    &=
    \int_0^\infty
      \paren[\Big]{\int_0^L\frac{\rho |g(t)|}{\rho^2+t^2}\,dt}^2\,d\rho .
  \end{align}
  The operator on $(0,\infty)$ with kernel
  \begin{align}
    K(\rho,t):=\frac{\rho}{\rho^2+t^2}
  \end{align}
  is bounded on $L^2(0,\infty)$.  Indeed, Schur's test with the weight $t^{-1/2}$ gives
  \begin{align}
    \int_0^\infty K(\rho,t)t^{-1/2}\,dt=C\rho^{-1/2},
    \qquad
    \int_0^\infty K(\rho,t)\rho^{-1/2}\,d\rho=Ct^{-1/2}.
  \end{align}
  Thus the contribution of $\Omega_-$ is bounded by $C\norm{g}_{L^2(S)}^2$, and the same estimate holds for $\Omega_+$.  Combining the strip and endpoint estimates proves \cref{eq:segment-usfe-appendix}.
\end{proof}

\end{document}
